% s12 Continuum subsequences and LIM states
\section{Continuum subsequences and LIM states}

The axioms of the preceding sections rest on the existence of the fixed-cube limits $\lim_k$: on \cref{hyp:s03:EC} in the curvature sector, on \cref{hyp:s03:W1} on the Wick tube. This section removes that dependence. The $k$-uniform bounds of \cref{lem:s05:C2} and \eqref{eq:s11:W} make the fixed-cube functionals normal families, so one subsequence of spacings serves every cube and every family at once; along it the cube limit $J\to\infty$ exists because the cone functionals are monotone in the cube at fixed $k$ (\cref{lem:s02:cut}). The limit state of such a subsequence (a \emph{LIM state}) satisfies every axiom proved above, with no hypothesis in the curvature sector and, on the Wick tube, under \Tthree\ and two further hypotheses; \cref{hyp:s03:EC} (resp.\ \cref{hyp:s03:W1}) adds exactly the identification with the iterated limit and the independence of the subsequence. The section is written on the comparison convention $\NDa$ (\cref{def:s02:acube}), in which the tube lemmas and the rotation part are written; King's convention is \cref{cor:s12:KK}. Notation is that of \cref{def:s02:sources,def:s02:tube,def:s11:cells} and \eqref{eq:s11:constants}: $a^{(a)}_{k,J}(z)=\ava{e^{\sum_jz_jF(f_j)}}$ for a family $f=(f_1,\dots,f_N)$ of real $2$-forms; $Z^a_{k,J}(\zeta)$ for $\cF=(f_1,\dots,f_n;h_1,\dots,h_p)$ and $\zeta=(z,w)\in\C^{n+p}$ in the scheme $a$; $g(\zeta)$, $k(\zeta)=k_+-k_-$, $\amp(\zeta)$; $U_J\defeq\{x:\dist_\infty(x,\R^3\setminus\CJ)>1\}$, $J_0(f)$ the least $J$ with all supports in $U_J$, and $R(\cF)$ the least $R$ with all supports in $Q_R=[-R,R]^3$. Two standard facts \cite{Hor90} are used by name: the \emph{identity theorem} (a holomorphic function on a connected open $D\subset\C^m$ vanishing on a non-empty relatively open subset of $D\cap\Sigma$, $\Sigma$ a maximal totally real affine subspace such as $\R^m$ or $(\mathrm i\R)^{m_1}\times\R^{m_2}$, vanishes on $D$) and \emph{Montel--Vitali} (a locally bounded family of holomorphic functions on $D$ is normal, and converges locally uniformly on $D$ if it converges on such a set).

\subsection{Definitions and hypotheses}

\begin{definition}[continuum subsequences]\label{def:s12:contsub}
An increasing sequence $\cK=(k_m)_{m\ge1}\subset\N$ of spacings is a \emph{continuum subsequence} (of the comparison family) if for every finite family $f$ of real $2$-forms and every $J\ge J_0(f)$ the sequence $(a^{(a)}_{k_m,J})_m$ converges locally uniformly on $\C^N$; the limit is written $Z^{\cK}_J(\sum_jz_jf_j,0)$ and depends on the family only through the total complex source, as each $a^{(a)}_{k,J}$ does. Fix a scheme $a$. $\cK$ is a \emph{continuum subsequence for the Wick sector} (Form~(B)) if it is a continuum subsequence and, for every finite family $\cF$ of real $2$-forms and scalars and every $J$ with $MJ>R(\cF)+\ell_c$, the sequence $(Z^a_{k_m,J})_m$ converges locally uniformly on the tube $\Dtube$ of $\cF$; the limit is $Z^{a,\cK}_J$. (The strict inequality replaces the threshold $J_1(R)$ of \cref{thm:s11:W} by at most one cube step; only tails in $J$ are used.) Form~(A), convergence on all of $\C^{n+p}$ for the same $J$, is the optional extension of \cref{thm:s23:FormA} and is not used here.
\end{definition}

\begin{definition}[LIM states]\label{def:s12:LIMstate}
The \emph{LIM state} of a continuum subsequence $\cK$ is $\ZLIM{\cK}\defeq\lim_{J\to\infty}Z^{\cK}_J$ on $\C^N$ (curvature sector), resp.\ $Z^{a,\cK}_\infty\defeq\lim_JZ^{a,\cK}_J$ on the tube (Wick sector); \cref{lem:s12:LIM1} gives existence. A subsequence of a continuum subsequence is a continuum subsequence with the same fixed-cube limits, hence the same LIM state.
\end{definition}

\begin{definition}[the limit state of the rotation theorem]\label{def:s12:o3LIM}
For a family $f$ and $J\ge J_0(f)$ let $L^f_J$ be the set of locally uniform limits on $\C^N$ of subsequences of $(a^{(a)}_{k,J})_k$. The rotation theorem (\cref{hyp:s03:RC}) gives an entire $Z^{\mathrm{LIM}}_\infty(\cdot,0)$ to which \emph{every} choice $\zeta_J\in L^f_J$ converges locally uniformly as $J\to\infty$, the subsequence being allowed to depend on $J$. With $J$-dependent subsequences one only has $\liminf_ka^{(a)}_{k,J}(\mathrm iu)\le\zeta_J(\mathrm iu)\le\limsup_ka^{(a)}_{k,J}(\mathrm iu)$, and the $J$-limit is pinned only if the $J$-limits of the two bounds agree, which is the content of that theorem. A LIM state in the sense of \cref{def:s12:LIMstate} uses one $\cK$ common to all $J$ and all families, along which the fixed-cube limits are monotone in $J$ (\cref{lem:s12:LIM1}), and needs no hypothesis; \cref{prop:s12:ID}(a) relates the two.
\end{definition}

The curvature sector needs no hypothesis. On the Wick tube the bound \eqref{eq:s11:W} of \cref{thm:s11:W} is used; it rests on \Tthree\ and the torus inputs (\cref{hyp:s03:T1,hyp:s03:T2,hyp:s03:T3}, which give \Itwo\ with amplitude $s_1\in(0,\sigma_1)$ and constant $C_U$ through \cref{prop:s10:U}; \Ithree\ is \cref{prop:s10:RP} and \Ifour\ is \cref{lem:s02:BFS}) and on one instance of \cref{hyp:s03:W1}, isolated as follows.

\begin{hypothesis}[\Bcell, the cell bound]\label{hyp:s12:Bcell}
In the scheme $a$, $B_{\mathrm{cell}}\defeq\sup_k\ln N^a_{k,C_{L_c}}(0,-2H'\chi_0)<\infty$, for the one cell $C_{L_c}$ (side $\ell_c\ge10$) and the one negative Wick source $-2H'\chi_0$ of \cref{def:s11:cells}.
\end{hypothesis}
This is the boundedness form of \cref{hyp:s03:W1} at one cube and one source, and the only place where it enters the proof of \eqref{eq:s11:W} (through \cref{lem:s11:W4,lem:s11:W5,prop:s11:W8}; $K_-$ is built from $B_{\mathrm{cell}}$). It follows from \cref{hyp:s03:W1} for an admissible $a$.

\begin{hypothesis}[\SCH, the scheme]\label{hyp:s12:SCH}
Either $a=z$, the reference scheme in which \cref{hyp:s03:T2} is computed, or $\sup_k|z_k-a_k|<\infty$; in the second case the constants of \eqref{eq:s11:W} hold with $C_U$ replaced by $C_U+s_1\sup_k|z_k-a_k|$.
\end{hypothesis}
For an admissible scheme (\cref{def:s03:adm}) \SCH\ holds; for LIM states with $a\ne z$ it is a genuine further hypothesis.

\begin{hypothesis}[\UNIQ, independence of the subsequence]\label{hyp:s12:UNIQ}
Any two continuum subsequences $\cK,\cK'$ have the same LIM state: $\ZLIM{\cK}(\cdot,0)=\ZLIM{\cK'}(\cdot,0)$ for every family.
\end{hypothesis}
\UNIQ\ follows from \cref{hyp:s03:RC} and from \cref{hyp:s03:EC} (\cref{prop:s12:ID}(a),(b)). Without them, every LIM state lies on $(\mathrm i\R)^N$ between $\lim_J\liminf_ka^{(a)}_{k,J}$ and $\lim_J\limsup_ka^{(a)}_{k,J}$ (both exist, by monotonicity in $J$), and \UNIQ\ holds if these coincide at every real argument, an entire function being determined on $(\mathrm i\R)^N$; the converse is not claimed. The axioms below hold for each LIM state separately and say nothing about independence of $\cK$, which is among the open problems.

\subsection{Existence of continuum subsequences and the cube limit}

\begin{lemma}[continuum subsequences exist]\label{lem:s12:LIM0}
\begin{enumerate}[label=\textup{(\alph*)},leftmargin=*,nosep]
\item \textup{(Curvature sector; no hypothesis.)} Every increasing sequence of spacings has a continuum subsequence. For every continuum subsequence $\cK$, family $f$ and $J\ge J_0(f)$: $Z^{\cK}_J(\cdot,0)$ is entire, $|Z^{\cK}_J(\sum_jz_jf_j,0)|\le e^{\norm{g(z)}_2^2}$, $Z^{\cK}_J(x,0)$ is real and $\ge1$ at real $x$, and $0\le Z^{\cK}_J(\mathrm it,0)\le1$ at imaginary $\mathrm it$. \emph{Several models at once:} for every countable set $\cM$ of lattice models of King's type (Neumann--Dirichlet cube states with a site-independent real bare scalar coefficient, photon mass $\mu>0$, $\lambda>0$; for instance the models at the photon masses of a countable set, with the counterterm fixed) one $\cK$ is a continuum subsequence for every member of $\cM$.
\item \textup{(Wick sector, Form~(B).)} Assume \Tthree, \Bcell\ and \SCH. Then every increasing sequence of spacings has a continuum subsequence for the Wick sector, and for every such $\cK$, family $\cF$ and $J$ with $MJ>R(\cF)+\ell_c$, $Z^{a,\cK}_J$ is holomorphic on $\Dtube$ and
\begin{equation}\label{eq:s12:Winf}
|Z^{a,\cK}_J(\zeta)|\le\exp\big(\norm{g(\zeta)}_2^2+\kappa\norm{k_+(\zeta)}_{L^1}+K_-\wnorm{k_-(\zeta)}\big)\qquad(\zeta\in\Dtube),
\end{equation}
with the constants of \eqref{eq:s11:constants}. \emph{Several models at once:} for a countable set of models all satisfying \eqref{eq:s11:W} with one set of constants, one $\cK$ serves all.
\end{enumerate}
\end{lemma}

Form~(A), which needs in addition the existence of the limits at every real coefficient vector (\Wtwo(a)), is \cref{thm:s23:FormA}. The subsequence is chosen in Step~2 below and nowhere else (\cref{rem:s12:K9}).

\begin{proof}
\emph{(a).} Inputs: \cref{lem:s05:C2} ($\sup_{k\ge0,\,J\ge J_0(f)}\sup_K|a^{(a)}_{k,J}|<\infty$ for every family and compact $K\subset\C^N$); \cref{lem:s05:C1} at second order ($\ava{F(u)^2}\le2\norm u^2_{2,k}$ for real $u$ with plaquettes in $\XJ$); charge conjugation $(A,\vartheta)\mapsto(-A,-\vartheta)$, a symmetry of every Neumann--Dirichlet state, so $\ava{F(u)}=0$ and $\ava{e^{\mathrm iF(u)}}$ is real.
\emph{Step 1 (a countable core).} Let $\cB$ be the countable set of finite unions of closed boxes with rational corners, for $B\in\cB$ let $S_0(B)$ be a countable $\norm\cdot_\infty$-dense subset of $C^\infty_c(\operatorname{int}B;\Lambda^2)$, and let $\cP$ be the countable set of pairs $(J,u')$, $u'$ a finite family from $\bigcup_BS_0(B)$, $J\ge J_0(u')$.
\emph{Step 2 (Montel and the diagonal).} For $(J,u')\in\cP$, $(a^{u'}_{k,J})_k$ is locally bounded on $\C^N$ uniformly in $k$ (\cref{lem:s05:C2}), hence normal; enumerate $\cP$, extract successively, and let $\cK=(k_m)$ be the diagonal sequence: $(a^{u'}_{k_m,J})_m$ converges locally uniformly for every $(J,u')\in\cP$.
\emph{Step 3 (every family, imaginary coefficients).} Let $f$ be a family, $J\ge J_0(f)$, and $B\in\cB$ with $\bigcup_j\supp f_j\subset\operatorname{int}B\subset B\subset U_{J_0(f)}$ (so every $u'\in S_0(B)$ is admissible at every $J\ge J_0(f)$). For real $2$-forms $u,u',g$ supported in $B$ and every $k$ with $\eps_k\le1$,
\[
|\ava{e^{\mathrm iF(u)}}-\ava{e^{\mathrm iF(u')}}|\le{\ava{F(u-u')^2}}^{1/2}\le\sqrt2\norm{u-u'}_{2,k},\qquad \norm g^2_{2,k}\le V_B\norm g_\infty^2,\quad V_B\defeq3\operatorname{vol}(N_{1/2}(B)),
\]
the first by $|e^{\mathrm ix}-1|\le|x|$ and the second-order bound, the second because the cubes of side $\eps_k$ about the centres of the plaquettes of one orientation on which $g\ne0$ are disjoint and lie in the closed $\ell^\infty$-neighbourhood $N_{1/2}(B)$. Given $t\in\R^N$ and $\eta>0$ pick $u'\in S_0(B)$ with $\norm{\sum_jt_jf_j-u'}_\infty\le\eta$; then $|a^f_{k_m,J}(\mathrm it)-a^f_{k_{m'},J}(\mathrm it)|\le2\sqrt{2V_B}\,\eta+|{\ava{e^{\mathrm iF(u')}}}_{k_m}-{\ava{e^{\mathrm iF(u')}}}_{k_{m'}}|$, and the last term tends to $0$ by Step~2. So $(a^f_{k_m,J}(\mathrm it))_m$ converges for every $t$.
\emph{Step 4 (all complex coefficients).} $(a^f_{k_m,J})_m$ is locally bounded on $\C^N$ and converges on the maximal totally real $(\mathrm i\R)^N$; Montel--Vitali gives locally uniform convergence to an entire $Z^{\cK}_J(\cdot,0)$. The bound $|a^{(a)}_{k,J}(z)|\le e^{\norm{g(z)}^2_{2,k}}$ (\cref{thm:s05:C}(a)) passes to the limit with $\norm g_{2,k}\to\norm g_2$ \eqref{eq:s02:riemann}; realness at real points passes; $a^{(a)}_{k,J}(x)=\ava{e^{F(u_x)}}\ge e^{\ava{F(u_x)}}=1$ by Jensen, and $[1,\infty)$ is closed; $0\le a^{(a)}_{k,J}(\mathrm it)\le1$ by the current representation of \cref{lem:s02:GD}(a), $\ava{e^{\mathrm i\ip vA}}=e^{-\ip v{Cv}/2}\,\E^\sharp\cosh\ip{Cv}{\kappa}\ge0$, and $|e^{\mathrm iF}|=1$. Dependence only on the total source holds at every $(k,J)$, hence in the limit.
\emph{Step 5 (several models).} The constant of \cref{lem:s05:C2}, the bound $e^{\norm g^2_{2,k}}$ and the Lipschitz constant $\sqrt2$ are model-free: \cref{thm:s05:C}(a) is \cref{lem:s02:GD}(a) for a measure of the class of \cref{lem:s02:BFS} plus \cref{lem:s05:C1} for the Neumann--Dirichlet photon covariance, valid for every real site-independent bare coefficient (the matter weight does not enter) and every $\mu>0$. For countable $\cM$ run Step~2 over $\cM\times\cP$ and Steps~3--4 member by member.

\emph{(b).} Inputs: under \Itwo--\Ifour\ and $B_{\mathrm{cell}}<\infty$, for every family with supports in $Q_R$, every $k$, $J\ge J_1(R)$ and $\zeta\in\Dtube$, \eqref{eq:s11:W} gives $|Z^a_{k,J}(\zeta)|\le Z^a_{k,J}(\Re\zeta)\le\exp(\norm g^2_{2,k}+\kappa\norm{k_+}_{1,k}+K_-\wnorm{k_-})$, with \Itwo\ from \Tthree\ by \cref{prop:s10:U} and \SCH\ for the scheme; on compacts of the tube the right side is bounded uniformly in $k$ (\cref{lem:s02:E}).
\emph{Steps 1--2.} As in (a), with $S^{\sW}_0(B)$ a countable $\norm\cdot_\infty$-dense subset of $C^\infty_c(\operatorname{int}B;\Lambda^2)\times C^\infty_c(\operatorname{int}B)$ and $\cP^{\sW}$ the countable set of pairs $(J,\cF')$, $\cF'$ a finite family from $\bigcup_BS^{\sW}_0(B)$, $MJ>R(\cF')+\ell_c$. For such a pair $J\ge J_1(R)$ for every $R\in(R(\cF'),MJ-\ell_c]$, so \eqref{eq:s11:W} holds for $\cF'$ at every $k$ and $(Z^a_{k,J})_k$ is normal on the tube of $\cF'$; the diagonal over $\cP^{\sW}$, together with that of (a), gives $\cK$.
\emph{Step 3 (every family, imaginary coefficients).} Let $\cF=(f_i;h_j)$, $MJ>R(\cF)+\ell_c$, $R_J\defeq MJ-\ell_c>R(\cF)$, so that all supports lie in the open cube $Q^\circ_{R_J}$ and $J\ge J_1(R_J)$; choose $B\in\cB$ with the supports in $\operatorname{int}B\subset B\subset Q^\circ_{R_J}$, so that $(J,\cF')\in\cP^{\sW}$ for every $\cF'$ from $S^{\sW}_0(B)$. For $(s,s')\in\R^{n+p}$ put $\sigma\defeq(\sum_is_if_i,\sum_js'_jh_j)$, so that $Z^a_{k,J}(\mathrm is,\mathrm is')=\ava{e^{\mathrm i\Phi^a_k(\sigma)}}$ with $\Phi^a_k(g,h)\defeq F_{\Lk}(g)+W^a_{\Lk}(h)$. \emph{Lipschitz bound.} For a real pair $q=(g_q,h_q)$ supported in $Q_{R_J}$, $(t_1,t_2)\mapsto Z^a_{k,J}(t_1g_q,t_2h_q)$ is holomorphic, and on the polydisc $|t_1|\le\norm{g_q}_{2,k}^{-1}$, $|t_2|\le\sigma_0/\wnorm{h_q}$ one has $\norm{k_\pm}_\infty\le|\Re t_2|\,\norm{h_q}_\infty\le\sigma_0<H/2<\sigma_+/2$, so the polydisc lies in the tube and \eqref{eq:s11:W} gives $|Z^a_{k,J}|\le e^{1+K'_{\mathrm{lat}}\sigma_0}$ there, $K'_{\mathrm{lat}}\defeq\kappa C_0'+K_-$ ($\norm{k_+}_{1,k}\le C_0'\wnorm{k_+}$ by \cref{lem:s02:E}). The Taylor coefficients at $0$ are the second moments of $F(g_q)$ and $W^a(h_q)$, and Cauchy's estimates give
\[
{\ava{\Phi^a_k(q)^2}}^{1/2}\le C\,(\norm{g_q}_{2,k}+\wnorm{h_q}),\qquad C\defeq\sqrt2\,e^{(1+K'_{\mathrm{lat}}\sigma_0)/2}\max(1,\sigma_0^{-1}),
\]
uniformly in $k$, $J$ and $R_J$; hence $|Z^a_{k,J}(\mathrm i\sigma)-Z^a_{k,J}(\mathrm i\sigma')|\le C(\norm{g_{\sigma-\sigma'}}_{2,k}+\wnorm{h_{\sigma-\sigma'}})$ for real pairs supported in $Q_{R_J}$. \emph{Density.} Given $\eta>0$ pick $\sigma'\in S^{\sW}_0(B)$ with $\norm{g_\sigma-g_{\sigma'}}_\infty,\norm{h_\sigma-h_{\sigma'}}_\infty\le\eta$; with $V_B$ and $\wnorm h\le(1+3R_J^2)^2\norm h_\infty$ for $h$ supported in $Q_{R_J}$, $|Z^a_{k,J}(\mathrm i\sigma)-Z^a_{k,J}(\mathrm i\sigma')|\le C(\sqrt{V_B}+(1+3R_J^2)^2)\eta$ for every $k$, while $Z^a_{k_m,J}(\mathrm i\sigma')$, the core functional at the point $(\mathrm i,\mathrm i)$ of its tube, converges by Step~2. So $(Z^a_{k_m,J})_m$ converges on $(\mathrm i\R)^{n+p}$.
\emph{Step 4 (the tube).} $(Z^a_{k_m,J})_m$ is locally bounded on the connected open $\Dtube$ and converges on the maximal totally real $(\mathrm i\R)^{n+p}\subset\Dtube$; Montel--Vitali gives locally uniform convergence to a holomorphic $Z^{a,\cK}_J$, and \eqref{eq:s11:W} passes to the limit with $\norm g_{2,k}\to\norm g_2$, $\norm{k_+}_{1,k}\to\norm{k_+}_{L^1}$ (Riemann sums; $k_+$ is continuous with compact support) and sampled $\wnorm{k_-}\le$ continuum $\wnorm{k_-}$: this is \eqref{eq:s12:Winf}.
\emph{Step 5 (several models).} With one set of constants in \eqref{eq:s11:W} for all members, $C$ is model-free and Step~2 runs over the product of the model set with $\cP^{\sW}$.
\end{proof}

\begin{lemma}[the cube limit along $\cK$ is a full limit]\label{lem:s12:LIM1}
Let $\cK$ be a continuum subsequence. \textup{(Curvature sector.)} For every family, $Z^{\cK}_J(\cdot,0)\to\ZLIM{\cK}(\cdot,0)$ locally uniformly on $\C^N$ as $J\to\infty$; the limit is entire, $|\ZLIM{\cK}(\sum_jz_jf_j,0)|\le e^{\norm{g(z)}_2^2}$, real and $\ge1$ at real points, in $[0,1]$ at imaginary points; $J\mapsto Z^{\cK}_J(\mathrm it,0)$ is nondecreasing, so $Z^{\cK}_J\le\ZLIM{\cK}$ on $(\mathrm i\R)^N$. \textup{(Wick sector, Form~(B).)} Under the hypotheses of \cref{lem:s12:LIM0}(b), $Z^{a,\cK}_J\to Z^{a,\cK}_\infty$ locally uniformly on $\Dtube$; the limit is holomorphic, obeys \eqref{eq:s12:Winf}, is $>0$ at the real points of the tube, and $Z^{a,\cK}_J\le Z^{a,\cK}_\infty$ at the cone points ($z\in(\mathrm i\R)^n$, $w$ real, $k(\zeta)\ge0$).
\end{lemma}

\begin{proof}
\emph{Curvature sector.} For $t\in\R^N$, $J\ge J_0(f)$ and every $k$: $a^{(a)}_{k,J}(\mathrm it)\le a^{(a)}_{k,J+1}(\mathrm it)\le1$ (first inequality of \cref{lem:s02:cut} with $Y=C_{J+1}\cap\Lk$, $X=\XJ$; the cut-out factor is the Neumann--Dirichlet state on $\CJ$ because $\dm_\eps$ is site-independent; realness by charge conjugation). Both sides converge along $\cK$, so $J\mapsto Z^{\cK}_J(\mathrm it,0)$ is nondecreasing, bounded by $1$, convergent. $\{Z^{\cK}_J\}_J$ is locally bounded on $\C^N$ by the $J$-free bound $e^{\norm{g(z)}_2^2}$ and converges on $(\mathrm i\R)^N$; Montel--Vitali gives the entire limit with the same bound; realness, $\ge1$ at real points and $[0,1]$ at imaginary points pass to the limit.

\emph{Wick sector.} (1) \emph{Split family.} Write $h_j=h_j^+-h_j^-$, $h_j^-\defeq\norm{h_j}_\infty\chi_j$, $h^+_j\defeq h_j+h^-_j\ge0$, with $\chi_j\in C^\infty_c$, $0\le\chi_j\le1$, $\chi_j=1$ on $\supp h_j$, $\supp\chi_j\subset N_{1/2}(\supp h_j)$; put $F_{k,J}(z,w^+,w^-)\defeq Z^a_{k,J}(\sum_iz_if_i,\sum_j(w^+_jh^+_j+w^-_jh^-_j))$ and $\cD''\defeq\{|\Re w^\pm_j|<\delta_j\}$ with $\sum_j\delta_j(\norm{h^+_j}_\infty+\norm{h^-_j}_\infty)<H/2$, a convex subset of the split family's tube. By \cref{lem:s12:LIM0}(b) for the split family, $F_{k_m,J}\to F^{\cK}_J$ locally uniformly on $\cD''$ for $MJ>R(\cF)+\frac12+\ell_c$, with \eqref{eq:s12:Winf}. (2) On $\Omega_+\defeq\cD''\cap((\mathrm i\R)^n\times(0,\infty)^{2p})$ the observable has imaginary curvature coefficients and nonnegative total scalar source, so $J\mapsto F_{k,J}$ is nondecreasing at every $k$ (first inequality of \cref{lem:s02:cut} with \cref{lem:s02:alg}(a), the scheme multiplier being common); along $\cK$, $J\mapsto F^{\cK}_J$ is nondecreasing and bounded on $\Omega_+$, hence convergent. (3) $\Omega_+$ is a non-empty open subset of the maximal totally real $(\mathrm i\R)^n\times\R^{2p}$ and $\cD''$ is connected: $F^{\cK}_J\to F^{\cK}_\infty$ locally uniformly on $\cD''$. (4) At $w^+=w$, $w^-=-w$ the lattice functionals equal $Z^a_{k,J}(\zeta)$ exactly (equal total sources), so $Z^{a,\cK}_J$ converges locally uniformly on $\cD'\defeq\{|\Re w_j|<\delta_j\}\times\C^n\subset\Dtube$; $\{Z^{a,\cK}_J\}_J$ is locally bounded on $\Dtube$ by the $J$-free \eqref{eq:s12:Winf} and converges on the open $\cD'$; Montel--Vitali on the connected tube gives $Z^{a,\cK}_J\to Z^{a,\cK}_\infty$, holomorphic, with \eqref{eq:s12:Winf}. (5) \emph{Positivity at real tube points.} For real $\zeta\in\Dtube$ pick $\epsilon>0$ with $-\epsilon\zeta\in\Dtube$ and let $C_\epsilon\defeq\sup_k(\norm g^2_{2,k}+\kappa\norm{k_+}_{1,k}+K_-\wnorm{k_-})|_{-\epsilon\zeta}<\infty$, the exponent of \eqref{eq:s11:W} at $-\epsilon\zeta$ (bounded in $k$ by \cref{lem:s02:E}, free of $J$). At every $k$, $t\mapsto\ln Z^a_{k,J}(t\zeta)$ is convex on $\R$ (H\"older, \cref{lem:s02:alg}(d)) and vanishes at $t=0$, so $1\le Z^a_{k,J}(\zeta)^{\epsilon/(1+\epsilon)}Z^a_{k,J}(-\epsilon\zeta)^{1/(1+\epsilon)}\le Z^a_{k,J}(\zeta)^{\epsilon/(1+\epsilon)}e^{C_\epsilon/(1+\epsilon)}$, i.e.\ $Z^a_{k,J}(\zeta)\ge e^{-C_\epsilon/\epsilon}$ for every $k$ and $J$; the bound passes along $\cK$ and then in $J$: $Z^{a,\cK}_\infty(\zeta)\ge e^{-C_\epsilon/\epsilon}>0$. (The product inequality of \cref{lem:s02:alg}(e) is not used: the tube is asymmetric, $H=\sigma_+/16$, and for real $\xi\in\Dtube$ with $\norm{k_+(\xi)}_\infty\ge H/2$ the point $-\xi$ lies outside it, where \cref{def:s12:contsub} gives no limit.) (6) At a cone point of the original family $J\mapsto Z^a_{k,J}(\zeta)$ is nondecreasing at every $k$ (first inequality of \cref{lem:s02:cut} with \cref{lem:s02:alg}(a)); pass along $\cK$.
\end{proof}

\subsection{The axioms for LIM states}

\begin{theorem}[LIM states of the comparison family at fixed photon mass]\label{thm:s12:LIM}
Let $e>0$, $\lambda,m,\lmu>0$, $\eps_k=5^{-k}$, $\CJ=[-MJ,MJ)^3$, and the states $\NDa_{k,J}$ of \cref{def:s02:acube}.

\textbf{Curvature sector (no hypothesis).} Every increasing sequence of spacings has a continuum subsequence. For every continuum subsequence $\cK$ and every finite family of real $2$-forms,
\[
\ZLIM{\cK}\Big(\sum_jz_jf_j,0\Big)=\lim_{J\to\infty}\lim_{m\to\infty}\Big\langle e^{\sum_jz_jF_{\Lambda_{k_m}}(f_j)}\Big\rangle^{(a)}_{k_m,J}
\]
exists locally uniformly in $z\in\C^N$, is entire, and depends on the family only through the total source. Its Schwinger functions $S^{\cK}_n(f_1,\dots,f_n)\defeq\partial_{z_1}\cdots\partial_{z_n}\ZLIM{\cK}(\sum_jz_jf_j,0)|_{z=0}$ satisfy:
\begin{enumerate}[label=\textup{(\roman*)},leftmargin=*,nosep]
\item \textup{(\OSz.)} $|\ZLIM{\cK}(\sum_jz_jf_j,0)|\le e^{\norm{g(z)}_2^2}$ (no dependence on $\lmu,e,\lambda,m$) and $\le\exp(\half\lmu^{-2}\norm{\sum_j(\Re z_j)\dstar f_j}_2^2)$; $|S^{\cK}_n|\le(\sqrt2e)^n(n!)^{1/2}\prod_j\norm{f_j}_2\le(e\pi/2)^n(n!)^{1/2}\prod_j\fournorm{f_j}$; the $S^{\cK}_n$ are tempered distributions, and \Ezz\ holds with $s=4$ and every $\beta>\half$.
\item \textup{(\OSo.)} $S^{\cK}_n$ is invariant under all translations and under $\Oh$ about every point, $F$ transforming as a $2$-form.
\item \textup{(\OSt.)} Reflection positivity about every coordinate plane $\{x_\nu=c\}$: for exponentials $E_a=e^{\sum_lz_{al}F(f_{al})}$ with arbitrary complex $z_{al}$ and $f_{al}$ supported in $\{x_\nu>c\}$, $(S(\overline{\theta_cE_a}E_b))_{ab}$ is positive semidefinite, and $S(\theta_c(P)P)\ge0$ for polynomials $P$ in such fields.
\item \textup{(\OSth.)} $S^{\cK}_n$ is real, symmetric, and zero for odd $n$.
\item \textup{(The measure.)} For every even $r\ge14$ there is exactly one Borel probability measure $\mu^{\cK}_F$ on $\cH^{-r}(\R^3;\Lambda^2)$ whose moments on $C^\infty_c$ are the $S^{\cK}_n$; its characteristic functional is $\ZLIM{\cK}(\mathrm i\,\cdot,0)$ extended by continuity to $\cS$, and its image on $H^{-r}_{\mathrm{loc}}$ is the unique measure there with this characteristic functional and the unique one with these moments. For every $\chi\in C^\infty_c$ and box $Q$ containing $\supp\chi$ in its interior, the laws of $\chi F_{\Lambda_{k_m}}$ under $\NDa_{k_m,J}$ converge weakly on $H^{-r}_{\mathrm{per}}(Q;\Lambda^2)$, \emph{first as $m\to\infty$ along $\cK$} at every fixed $J\ge J_0(\chi)$ and then as $J\to\infty$, to the law of $\chi\Phi$ under $\mu^{\cK}_F$.
\item \textup{(Weak \OSf.)} $\mu^{\cK}_F$ is $\Cub$-invariant and satisfies the clauses of \cref{thm:s09:S2K}(b) with $\mu^{\cK}_F$ in place of $\mu^K_S$: it is an integral of translation-ergodic, reflection-positive, $\theta_j$-invariant components clustering in the mean and under diagonal separation, with a unique joint vacuum per coordinate direction and the component growth bound; almost every component satisfies \OSf\ in the weak form; the $\Oh$-invariant decomposition into mixtures of $k\in\{1,2,3,6\}$ ergodic measures exists.
\item \textup{(Euclidean clause.)} $\ZLIM{\cK}$ does not depend on $\cK$, equals the limit state $Z^{\mathrm{LIM}}_\infty$ of \cref{def:s12:o3LIM}, is $\Euc$-invariant for all complex coefficients (\cref{hyp:s03:RC}), and the Euclidean forms of (ii) (full $\Euc$), (iii) (reflection positivity about every plane) and (vi) (almost every component $\Euc$-invariant, reflection positive about every plane, clustering in any manner, with a unique vacuum in every direction; the vanishing of a component's odd moments is not claimed) hold.
\end{enumerate}

\textbf{Wick sector, Form~(B).} Assume \Tthree, \Bcell\ and \SCH. Every increasing sequence of spacings has a continuum subsequence for the Wick sector, and for every such $\cK$ and family the LIM state $Z^{a,\cK}_\infty$ exists locally uniformly on $\Dtube$, is holomorphic there, obeys \eqref{eq:s12:Winf}, and (i)--(vi) hold for the joint Schwinger functions $S^{\cK}_{n,p}(f;h)=\partial^{(n)}_z\partial^{(p)}_wZ^{a,\cK}_\infty(0)$ with: in (i) the mixed bounds $|S^{\cK}_{n,p}|\le e^{K'\sigma_0}(\sqrt2e)^n(n!)^{1/2}(e/\sigma_0)^p\,p!\prod_i\norm{f_i}_2\prod_j\wnorm{h_j}$, $\sigma_0=H/4$, $K'=\kappa\pi^2+K_-$, and \Ezz\ with $s=4$ and every $\beta>1$; in (ii) $W^a$ a scalar, the identity holding on the tube; in (iii) exponentials with arbitrary complex curvature coefficients and Wick coefficient vectors $\zeta_1,\dots,\zeta_N$ of \emph{amplitude $\amp(\zeta_i)<H/2$} (in particular all purely imaginary Wick coefficients), and the polynomial form; in (iv) $S^{\cK}_{n,p}=0$ for odd $n$; in (v) the measures on $\cH^{-r}(\R^3;\Lambda^2\oplus\R)$ and $H^{-r}_{\mathrm{loc}}$, the localised lattice fields $\chi(F_{\Lk},W^a_{\Lk})$ converging along $\cK$ and then in $J$; (vi) unchanged; and \emph{no Euclidean clause}. \emph{Schemes:} if $a,a'$ both satisfy the hypotheses, every continuum subsequence has a further subsequence $\cK'$ along which $a'_k-a_k\to\kappa^{\cK'}$, and $Z^{a',\cK'}_\infty(\zeta)=e^{-\kappa^{\cK'}\sum_jw_j\int h_j}Z^{a,\cK'}_\infty(\zeta)$: the two LIM states differ by the shift $W\mapsto W-\kappa^{\cK'}\int(\cdot)$ of the Wick component, and joint cumulants of order $\ge2$ are scheme-free.
\end{theorem}

\begin{proof}
Each clause is the corresponding proof for the iterated limit (\cref{thm:s05:C,prop:s06:translations,thm:s06:S1K,thm:s08:OS2C,thm:s08:OS2WT,prop:s08:OS3,thm:s09:S2K}) with exactly one of three replacements: ($\alpha$) a fixed-$(k,J)$ identity or inequality between Neumann--Dirichlet functionals of explicit finite families is passed along $\cK$ (both sides converge by \cref{lem:s12:LIM0}, which covers \emph{every} family); ($\beta$) Montel on a $k$-uniform bound, the identification of the limit by convergence at real points being replaced by the choice of $\cK$; ($\gamma$) a step about the limit object only, unchanged and not repeated here. The only uniqueness used is that of the monotone cube limit along $\cK$ (\cref{lem:s12:LIM1}).

\emph{(i).} \cref{lem:s12:LIM0,lem:s12:LIM1} give existence, entireness and the first bound; the second, $\exp(\half\lmu^{-2}\norm{\dstar_kg}^2_{2,k})$ at every $(k,J)$ (\cref{thm:s05:C}(a)), passes along $\cK$ and $J$ with $\norm{\dstar_kg}_{2,k}\to\norm{\dstar g}_2$. The Schwinger bounds, the extension to $\cS$, the kernel theorem and \Ezz\ are the arguments of \cref{thm:s05:C} (Wick sector: \cref{thm:s11:W}(ii), on the polydisc inside the tube) run on the limit function with its bound ($\gamma$).

\emph{(ii), translations.} For $b\in5^{-k_0}\Z^3$, $|b|_\infty\le M$, $k\ge k_0$: $Z^a_{k,\CJ}(\tau_b\cdot\zeta)=Z^a_{k,\CJ-b}(\zeta)$ (exact lattice covariance) and $C_{J-1}\cap\Lk\subset(\CJ-b)\cap\Lk\subset C_{J+1}\cap\Lk$, so for a cone point $s$ supported in $U_{J_0}$ and $J\ge J_0+1$, $Z^a_{k,C_{J-1}}(s)\le Z^a_{k,\CJ}(\tau_b\cdot s)\le Z^a_{k,C_{J+1}}(s)$ (first inequality of \cref{lem:s02:cut}, twice). The three terms are functionals of the families $s$ and $\tau_b\cdot s$ (supports in $Q_{R+M}$, the same tube), so they converge along $\cK$ ($\alpha$), and \cref{lem:s12:LIM1} sends both ends to $\ZLIM{\cK}(s)$: $\ZLIM{\cK}(\tau_b\cdot s)=\ZLIM{\cK}(s)$ at cone points, with no uniqueness of the $k$-limit. All coefficients ($\gamma$): the difference of the two limit functions is entire and vanishes on $(\mathrm i\R)^N$ (curvature), resp.\ for the split family is holomorphic on the connected $\cD''$ and vanishes on $\Omega_+$, so the identity theorem, the diagonal (4) of \cref{lem:s12:LIM1} and the identity theorem on the connected tube give it on $\Dtube$. All $b\in\bigcup_k5^{-k}\Z^3$ by iteration; all $b\in\R^3$ at imaginary coefficients because the Lipschitz bounds of \cref{lem:s12:LIM0}, Step~3, pass to the limits ($|\ZLIM{\cK}(\mathrm iu,0)-\ZLIM{\cK}(\mathrm iu',0)|\le\sqrt2\norm{u-u'}_2$; $|Z^{a,\cK}_\infty(\mathrm i\sigma)-Z^{a,\cK}_\infty(\mathrm i\sigma')|\le C(\norm{g_{\sigma-\sigma'}}_2+\wnorm{h_{\sigma-\sigma'}})$) and $\tau_{b_n}s\to\tau_{b_0}s$ in $L^2$ and in $\wnorm\cdot$ for $s\in C^\infty_c$; complex coefficients by the identity theorem.

\emph{(ii), the cubic group.} Coordinate permutations are exact at every $(k,J)$ (\cref{lem:s02:pass}); pass along $\cK$ and $J$ ($\alpha$). For the reflection $\sigma_3$ and a cone point $(f,g)$ supported in $U_{J-1}$, at every $k$, $Z^a_{k,C_{J-1}}(\mathrm if,g)\le Z^a_{k,\CJ}(\mathrm i\sigma_3\cdot f,\sigma_3\cdot g)\le Z^a_{k,C_{J+1}}(\mathrm if,g)$ (the frame flip of \cref{lem:s02:pass} with the one-site shift $X_+=(\CJ+\eps_ke_3)\cap\Lk$ and two deletions in the cone direction, \cref{lem:s02:BFS}(a),(b)); the middle term is the functional of the reflected family (same tube), so ($\alpha$) and \cref{lem:s12:LIM1} give $\ZLIM{\cK}(\mathrm i\sigma_3\cdot f,\sigma_3\cdot g)=\ZLIM{\cK}(\mathrm if,g)$ on the cone (all of it at $g=0$; in Form~(B) at cone points with $\norm g_\infty<\sigma_+/2$): the reflections are not exact lattice symmetries; the sandwich replaces exactness and survives along $\cK$. Conjugation $\sigma_\nu=P\sigma_3P^{-1}$, group closure on the cone, the identity theorem off the cone (on $\C^N$ with $(\mathrm i\R)^N$; on $\cD''$ with $\Omega_+$), invariance about every point by the translations, and differentiation are ($\gamma$), as in \cref{thm:s06:S1K}.

\emph{(iii), reflection positivity.} Fix $\nu$, a mid-plane $c\in D$ (\cref{def:s07:reflection}) and the enlarged family $\cF$ of \cref{def:s08:enlarged} with the modified-box functionals $\Phi_k(\zeta)=Z^{a,(3.38)}_{k,\Omega_{J,k}}(\zeta)$, $\zeta\in\C^m$, and support radius $R_\cF$. The fixed-$k$ steps (reflection positivity of the modified box, \cref{lem:s07:RPbox}; the sandwich lemmas \cref{lem:s07:L,lem:s07:U}) concern the photon form and the inequalities of \cref{lem:s02:BFS} only and are unchanged. \emph{Curvature sector.} At $h=0$, $|\Phi_k(z)|\le e^{\norm{f_{\Re z}}^2_{2,k}}$ by \cref{lem:s02:GD}(a) and the modified photon form, uniformly in $k$ and $J$. ($\beta$) Extract from $\cK$ a further $\cK'$ along which $\Phi_k\to\Phi^J$ locally uniformly on $\C^m$; positive semidefiniteness of the Gram matrices at every $k\ge k_0(c,\delta)$ passes to $\Phi^J$. ($\alpha$) On $U$ (imaginary curvature coefficients) $a^{(a)}_{k,\CJ}(\zeta)\le\Phi_k(\zeta)\le a^{(a)}_{k,C_{J+2n}}(\zeta)$ at every $k$ (\cref{lem:s07:L,lem:s07:U}, cone halves); along $\cK'$ the ends converge to $Z^{\cK}_J,Z^{\cK}_{J+2n}$, so $Z^{\cK}_J\le\Phi^J\le Z^{\cK}_{J+2n}$ on $U$, and \cref{lem:s12:LIM1} sends both ends to $\ZLIM{\cK}$. $\{\Phi^J\}_J$ is locally bounded on $\C^m$ ($J$-free bound) and $U$ is a set of uniqueness, so $\Phi^J\to\ZLIM{\cK}$ locally uniformly, in particular at the matrix points, where $\Phi^J(\zeta^*_{ij})\to\ZLIM{\cK}(\theta\zeta^*_i\sqcup\zeta_j)$ (dependence through the total source only); positivity passes. $\cK'$ may depend on $J$; this is harmless, since the sandwich ends converge along all of $\cK$ and the bound on $\Phi^J$ is $J$-free, so neither $\lim_J\Phi^J$ nor the positivity depends on the choice. Every plane: translate by $(c'-c)e_\nu$ and use (ii). The polynomial form: Weierstrass along $\cK'$ and $J$ (\cref{prop:s08:poly}). \emph{Wick sector, amplitude $<H/2$, without \cref{hyp:s03:W1}.} Let $\zeta_1,\dots,\zeta_N$ be families of real sources supported in $\{x_\nu>c\}$ with $\amp(\zeta_i)<H/2$; the $(i,j)$ Gram entry is $\Phi_k(\zeta^*_{ij})$, with total scalar source $\sum_l\bar w_{il}\,h_{il}\circ\rho+\sum_lw_{jl}h_{jl}$. The bound of \cref{lem:s08:claim4}, which depends on \cref{hyp:s03:W1}, is replaced by \cref{lem:s12:claim4prime} below, after which the argument is completed.

\emph{(iv).} Symmetry: Taylor coefficients of a function of the total source. Realness: real at real points at every $(k,J)$, closed under limits (Form~(B): at the real points of the tube). Charge conjugation gives $Z^a_{k,J}(-z,w)=Z^a_{k,J}(z,w)$ at every $(k,J)$, hence for $Z^{\cK}_J$ and $\ZLIM{\cK}$ (the tube is invariant under $z\mapsto-z$); odd $z$-derivatives at $0$ vanish.

\emph{(v), (vi).} The moment bounds at $(k,J)$ (\cref{thm:s05:C}(a), resp.\ \eqref{eq:s11:W}, with Cauchy's estimates) are uniform in $k$ and $J$ above the thresholds, and the lattice Schwinger functions converge along $\cK$ then $J$ (Weierstrass). The measure is built from the limit functional ($\gamma$): $\cZ^{\cK}(f)\defeq\ZLIM{\cK}(\mathrm if)$ on real $f\in C^\infty_c$ is normalised, positive definite (at every $(k,J)$ the matrix $(Z^a_{k,J}(\mathrm i(f_m-f_l)))_{lm}$ is a Gram matrix; the entries converge along $\cK$ then $J$) and Lipschitz for $\sqrt2\norm\cdot_2$, resp.\ $C(\norm\cdot_2+\wnorm\cdot)$, hence continuous on $\cS$; Bochner--Minlos, moments by Lukacs \cite{Luk70} and uniform integrability, the Hermite carrier on $\cH^{-r}$, $r\ge14$ even, and the two uniqueness statements are the steps of \cref{thm:s09:S2K}(a) with \cite[lem:D:hermite(c), lem:D:realize, prop:discharge]{SymPhase}. \emph{Local laws along $\cK$} (($\beta$); the one weakened clause): tightness of the laws of $\chi\Phi^a_{k,J}$ on $H^{-r}_{\mathrm{per}}(Q;\cV)$ is uniform in $k$ and $J\ge J(\chi)$ from the second-moment bounds (Prokhorov \cite{Bil99}); the characteristic functionals $Z^a_{k,J}(\mathrm i\chi\psi)$ converge \emph{along $\cK$} to $Z^{a,\cK}_J(\mathrm i\chi\psi)$; a probability measure on $H^{-r}_{\mathrm{per}}$ is determined by its characteristic functional on trigonometric polynomials; so the laws converge along $\cK$, then in $J$ (\cref{lem:s12:LIM1}), to the law of $\chi\Phi$ under the measure just built (the periodisation identity of \cref{thm:s09:S2K}). (vi): $\Cub$-invariance of the measure is (ii) at imaginary coefficients plus uniqueness of the characteristic functional; (H3) of \cite{SymPhase} is (iii) at purely imaginary Wick coefficients ($\amp=0<H/2$) with the density of the exponential algebra \cite[lem:A:mc]{SymPhase}; then \cite[thm:D:decomp, thm:D:mean, thm:D:diag, prop:D:recon, cor:D:schwinger(iv), add:Ohdecomp, add:OS4(iii)]{SymPhase}, stated for a measure satisfying (H1)--(H3) and matched as in \cref{rem:s09:v1match}.

\emph{Schemes.} $\sup_k|a'_k-a_k|\le\sup_k|a'_k-z_k|+\sup_k|z_k-a_k|<\infty$ by \SCH\ for both. For a continuum subsequence $\cK$ for $a$, Bolzano--Weierstrass gives $\cK'\subset\cK$ with $a'_k-a_k\to\kappa^{\cK'}\in\R$. At every $(k,J)$, $Z^{a'}_{k,J}(\zeta)=e^{-\sum_jw_j(a'_k-a_k)\sS_{\Lk}(h_j)}Z^a_{k,J}(\zeta)$ (\cref{rem:s02:scheme}) with $\sS_{\Lk}(h_j)\to\int h_j$; so along $\cK'$ the $a'$-functionals converge wherever the $a$-functionals do, $\cK'$ is a continuum subsequence for $a'$ with $Z^{a',\cK'}_J=e^{-\kappa^{\cK'}\sum_jw_j\int h_j}Z^{a,\cK'}_J$, and then $J\to\infty$. (For admissible $a,a'$, $a'_k-a_k$ converges along the full sequence and $\kappa^{\cK'}=\kappa$.) The cumulant statement is that of \cref{thm:s13:main}. \Bcell\ transfers between schemes with bounded difference: $B^{a'}_{\mathrm{cell}}\le B^a_{\mathrm{cell}}+\sup_k|a'_k-a_k|\,\norm{2H'\chi_0}_{1,k}$.

\emph{(vii).} By \cref{prop:s12:ID}(a), $\ZLIM{\cK}=Z^{\mathrm{LIM}}_\infty$, which is $\Euc$-invariant for all $z\in\C^N$ by \cref{hyp:s03:RC}. (ii) in Euclidean form follows by differentiation; for (iii), a plane $P$ is $E\{x_\nu=c\}$ for a Euclidean motion $E$, $\theta_P=E\theta_{\nu,c}E^{-1}$, and $\ZLIM{\cK}(\theta_P\zeta^*_i\sqcup\zeta_j)=\ZLIM{\cK}(\theta_{\nu,c}(E^{-1}\zeta_i)^*\sqcup E^{-1}\zeta_j)$ by invariance at complex coefficients, the $E^{-1}\zeta_i$ being supported in $\{x_\nu>c\}$; (vi) is \cite[p:cor:rotcor]{SymPhase} with $\cZ^{\cK}$ $O(3)$-invariant, as in \cref{thm:s09:S2K}(c), with its proviso on the odd moments of a component.
\end{proof}

\begin{lemma}[a $k$- and $J$-uniform bound for the modified-box functionals on the tube of the enlarged family]\label{lem:s12:claim4prime}
Under \Tthree, \Bcell, \SCH, for every $k$, every $J\ge\max(J_0,J_1(R_\cF))$ and every $\zeta\in\cD^{\cF}_{\sigma_+,H}$ (the tube of $\cF$ in the sense of \cref{def:s02:tube}),
\[
|\Phi_k(\zeta)|\le\exp\big(2\norm{g(\zeta)}^2_{2,k}+\kappa\norm{k_+(\zeta)}_{1,k}+K_-\wnorm{k_-(\zeta)}\big).
\]
\end{lemma}
The curvature constant is $2$, not the $1$ of \eqref{eq:s11:W}, because the Cauchy--Schwarz split below doubles the curvature source and \cref{thm:s05:C}(a) is quadratic in it; nothing depends on its value.

\begin{proof}
$|\Phi_k(\zeta)|\le\Phi_k(\xi)$, $\xi\defeq\Re\zeta$ (\cref{lem:s02:alg}(d)). Let $k(\xi)$ be the total real scalar source at $\xi$ and $k_1\defeq k(\xi)_+$, $k_2\defeq k(\xi)_-$ its \emph{pointwise} parts, sampled on the lattice. Cauchy--Schwarz at the lattice level (\cref{lem:s02:alg}(f); the multiplier $e^{-a_k\sS(k_1-k_2)}$ splits by linearity of $\sS$): $\Phi_k(\xi)\le Z^{a,(3.38)}(2f_\xi,-2k_2)^{1/2}Z^{a,(3.38)}(0,2k_1)^{1/2}$. \emph{First factor} (class $\cD$): the $\cD$-half of \cref{lem:s07:L}, which needs only that the sources vanish off $\XJ$ ($J\ge J_0$), gives $\le Z^a_{k,\CJ}(2f_\xi,-2k_2)$; \cref{thm:s05:C}(a) with the real scalar source $-2k_2$ gives $\le e^{\norm{2f_\xi}^2_{2,k}}N^a_{k,\CJ}(0,-2k_2)$ (the same scheme multiplier on both sides); \cref{prop:s11:W8} with $q\defeq2k_2\ge0$, supported in $Q_{R_\cF}$, $J\ge J_1(R_\cF)$ and $\sup q=2\norm{k(\xi)_-}_\infty<H$ (this is $\zeta\in\cD^\cF_{\sigma_+,H}$; the proposition allows $\sup q\le H$) gives $\ln N^a_{k,\CJ}(0,-2k_2)\le K_-\wnorm{2k_2}$. Hence $\half\ln(\text{first factor})\le2\norm{g(\zeta)}^2_{2,k}+K_-\wnorm{k_-(\zeta)}$, as $f_\xi=g(\zeta)$ and $k_2=k_-(\zeta)$. \emph{Second factor} (class $\cI$): the $\cI$-half of \cref{lem:s07:U} with $B=C_{J+2n}\cap\Lk$ gives $\le Z^a_{k,C_{J+2n}}(0,2k_1)$, and \cref{cor:s11:W2prime}, uniform in $k$ and in the cube, with $0\le2k_1\le s_1$ on the lattice ($\norm{k_+(\zeta)}_\infty<\sigma_+/2$, sampled $\le$ continuum) gives $\ln Z^a_{k,C_{J+2n}}(0,2k_1)\le\kappa\norm{2k_1}_{1,k}$, so $\half\ln(\text{second factor})\le\kappa\norm{k_+(\zeta)}_{1,k}$. Add. On compacts of the tube the right side is bounded uniformly in $k$ (\cref{lem:s02:E}).
\end{proof}

The lemma is \eqref{eq:s11:W} (with curvature constant $2$) for the modified-box functional, transported through \cref{lem:s07:L,lem:s07:U}; the two factors live on $\CJ$ and $C_{J+2n}$, which is harmless as \cref{prop:s11:W8} and \cref{cor:s11:W2prime} are uniform in the cube, and the amplitude $H/2$ is exactly the tube condition $\norm{k(\xi)_-}_\infty\le H/2$ on the total source. \emph{The matrix points lie in the tube:} at $\xi^*_{ij}=\Re\zeta^*_{ij}$ the two sums of the total scalar source have disjoint supports, so $\norm{k(\xi^*)_\pm}_\infty\le\max(\amp(\zeta_i),\amp(\zeta_j))<H/2\le\sigma_+/2$, and by continuity a neighbourhood of the finitely many matrix points lies in $\cD^\cF_{\sigma_+,H}$. \emph{End of the proof of (iii), Wick sector.} ($\beta$) By \cref{lem:s12:claim4prime}, $(\Phi_k)_k$ is locally bounded on the open connected $\cD^\cF_{\sigma_+,H}$; extract $\cK'\subset\cK$ with $\Phi_k\to\Phi^J$ locally uniformly there; Gram positivity passes ($\cK'$ may depend on $J$; harmless as in the curvature sector). ($\alpha$) Let $U$ be the set of coefficient vectors with purely imaginary curvature coefficients and positive real scalar coefficients on the nonnegative $\tilde h_j$; $U\cap\cD^\cF\ne\emptyset$ ($k_-=0$, $k_+$ small), and it is a relatively open subset of the maximal totally real $(\mathrm i\R)^{m_F}\times\R^{m_h}$, a set of uniqueness in $\cD^\cF$. On $U$, $Z^a_{k,\CJ}(\zeta)\le\Phi_k(\zeta)\le Z^a_{k,C_{J+2n}}(\zeta)$ at every $k$ (\cref{lem:s07:L,lem:s07:U}, $\cI$-halves); the ends are functionals of $\cF$ at points of its tube, converging along $\cK'$ (\cref{lem:s12:LIM0}(b), $J$ above the threshold of $\cF$) to $Z^{a,\cK}_J,Z^{a,\cK}_{J+2n}$, which tend to $Z^{a,\cK}_\infty$ (\cref{lem:s12:LIM1}, locally uniformly on the tube of $\cF$, in particular on $U\cap\cD^\cF$). $\{\Phi^J\}_J$ is locally bounded on $\cD^\cF$ ($J$-free), so $\Phi^J\to Z^{a,\cK}_\infty$ locally uniformly there, in particular at the matrix points, where the limit is $Z^{a,\cK}_\infty(\theta\zeta^*_i\sqcup\zeta_j)$; positivity passes. Every plane by (ii) (the amplitude is translation invariant); the polynomial form by Weierstrass.

\begin{remark}[amplitudes in $[H/2,\sigma_+/8)$: the route through the fixed-cube limits, for a further subsequence]\label{rem:s12:amplitudes}
The proof of \cref{lem:s08:claim4} bounds the first factor on the Wick tube $T_{\cF'}$ of \cref{def:s08:enlarged} by the vertex argument at the cube $C_{J_0}$, which needs the fixed-cube limits in their boundedness form only. For amplitudes in $[H/2,\sigma_+/8)$ the matrix points lie in $T_{\cF'}$ but in general outside $\cD^{\cF}_{\sigma_+,H}$, the only domain on which \cref{def:s12:contsub} gives convergence along $\cK$; what follows is therefore stated here only for a further subsequence. Assume the boundedness in $k$ of the functionals of the class-$\cD$ families $(2f_\xi,-2k_2(\xi))$ at the one cube $C_{J_0}$ (with the coefficient-sign split of \cref{def:s08:enlarged}). The vertex argument then gives local boundedness of $(Z^a_{k,J})_k$ and $(\Phi_k)_k$ on $T_{\cF'}$ (first factors at $C_{J_0}$ by that boundedness, second factors by \cref{cor:s11:W2prime}); Montel, with a diagonal over the countably many $J$, yields a further subsequence $\cK'\subset\cK$, depending on the family, along which $Z^a_{k,J}\to Z^{a,\cK'}_J$ and $\Phi_k\to\Phi^J$ locally uniformly on $T_{\cF'}$ (on $T_{\cF'}\cap\cD^\cF$ the first limit is $Z^{a,\cK}_J$). The sandwich step runs on $T_{\cF'}$ with the uniqueness set $U\cap T_{\cF'}$, the $J$-free local bound of \cref{prop:s08:T}(a) on $T_{\cF'}$ serving the Vitali step in $J$, and \cref{prop:s08:T}(b) along $\cK'$ for the remaining planes; so \OSt\ holds up to $\amp(\zeta_i)<\sigma_W=\sigma_+/8$ for the LIM state of such a $\cK'$ on $T_{\cF'}$. Whether the convergence holds along $\cK$ itself on $T_{\cF'}$ is not claimed here; under the hypotheses of Form~(A), \cref{thm:s23:FormA} gives \OSt\ along $\cK$ with no amplitude restriction. The remark is used nowhere else.
\end{remark}

\subsection{Identification, and King's printed state}

\begin{proposition}[identification]\label{prop:s12:ID}
\begin{enumerate}[label=\textup{(\alph*)},leftmargin=*,nosep]
\item For every continuum subsequence $\cK$ and every family, $\ZLIM{\cK}(\cdot,0)=Z^{\mathrm{LIM}}_\infty(\cdot,0)$, the limit state of the rotation theorem (\cref{hyp:s03:RC}); in particular all LIM states coincide, \UNIQ\ holds, and they are $\Euc$-invariant. \Cref{hyp:s03:RC} does \emph{not} give \cref{hyp:s03:EC}.
\item The following are equivalent: (i) the comparison family is convergent at fixed cubes (\cref{def:s03:conv}), that is the existence in $(0,\infty)$ of the fixed-cube limits of the comparison family at real coefficients; (ii) for every $J$ and family, $a^{(a)}_{k,J}\to Z_J$ locally uniformly on $\C^N$ along the full sequence; (iii) every increasing sequence of spacings is a continuum subsequence. Under them $Z^{\cK}_J=Z_J$ and $\ZLIM{\cK}=Z_\infty\defeq\lim_JZ_J$ for every $\cK$: every LIM state is the iterated limit of the comparison family.
\item Assume \Tthree\ and let $a$ be admissible. Then \SCH\ holds (\cref{def:s03:adm}) and \Bcell\ holds (\cref{hyp:s03:W1} for $a$ at the one cell and source), every increasing sequence of spacings is a continuum subsequence for the Wick sector, and for every $\cK$, $Z^{a,\cK}_\infty=\Ginf$ of \cref{thm:s11:W} on $\Dtube$, which equals $Z^a_{\infty,T}$ of \cref{prop:s08:T} on $T_{\cF'}\cap\Dtube$ (\cref{prop:s08:T} lives on the Wick tube $T_{\cF'}$, a different domain). Once \SCH\ and \Bcell\ are given, the argument uses \cref{hyp:s03:W1} for $a$ alone.
\end{enumerate}
\end{proposition}

\begin{proof}
(a) Let $\cK$ be a continuum subsequence, $f$ a family, $J\ge J_0(f)$. By \cref{def:s12:contsub}, $Z^{\cK}_J(\cdot,0)$ is the locally uniform limit of $(a^{(a)}_{k_m,J})_m$, so $Z^{\cK}_J(\cdot,0)\in L^f_J$. By \cref{hyp:s03:RC} every choice $\zeta_J\in L^f_J$, $J\ge J_0(f)$, converges locally uniformly to $Z^{\mathrm{LIM}}_\infty(\cdot,0)$; take $\zeta_J\defeq Z^{\cK}_J(\cdot,0)$, which by \cref{lem:s12:LIM1} converges to $\ZLIM{\cK}(\cdot,0)$, and limits are unique. Hence any two continuum subsequences have the same LIM state, and invariance is the second clause of \cref{hyp:s03:RC}. That theorem pins the $J$-limit of the subsequential limits, not the fixed-cube limits themselves, so it does not give \cref{hyp:s03:EC}.

(b) (i)$\Rightarrow$(ii): $(a^{(a)}_{k,J})_k$ is locally bounded on $\C^N$ (\cref{lem:s05:C2}) and converges at every real point; $\R^N$ is a set of uniqueness for entire functions, and Montel--Vitali gives locally uniform convergence of the whole sequence to an entire $Z_J$. (ii)$\Rightarrow$(iii): every subsequence converges to the same limit. (iii)$\Rightarrow$(i): the full sequence is itself an increasing sequence, hence a continuum subsequence, so $a^{(a)}_{k,J}$ converges locally uniformly, and at real $x$ the limit lies in $[1,\infty)\subset(0,\infty)$ (\cref{lem:s12:LIM0}, Step~4). The last clause is \cref{lem:s12:LIM1}.

(c) Admissibility gives $\sup_k|z_k-a_k|<\infty$, hence \SCH; \cref{hyp:s03:W1} for $a$ at the cell gives $B_{\mathrm{cell}}<\infty$; so \eqref{eq:s11:W} holds. $(Z^a_{k,J})_k$ is locally bounded on the tube by \eqref{eq:s11:W} and converges at the real points of the tube by \cref{hyp:s03:W1}, a non-empty open subset of $\R^{n+p}$ and a set of uniqueness in the connected tube; Montel--Vitali: the whole sequence converges locally uniformly on the tube (\cref{thm:s11:W}(i)), so every increasing sequence is a continuum subsequence for the Wick sector with $Z^{a,\cK}_J=G_J$, and $Z^{a,\cK}_\infty=\Ginf$ by \cref{lem:s12:LIM1}; on $T_{\cF'}\cap\Dtube$ this is the same $J$-limit of the same $G_J$ as $Z^a_{\infty,T}$ (\cref{prop:s08:T}(a)).
\end{proof}

\begin{corollary}[the iterated limits of the comparison family]\label{cor:s12:K}
Assume that the comparison family is convergent at fixed cubes (\cref{def:s03:conv}). Then every increasing sequence of spacings is a continuum subsequence, all LIM states coincide with the iterated limit $Z_\infty(f,0)=\lim_J\lim_k\ava{e^{F(f)}}$ of the comparison family, and the curvature sector (i)--(vi) of \cref{thm:s12:LIM} holds for it with the lattice laws of (v) converging along the \emph{full} sequence of spacings; (vii) holds as well. For general sources assume \Tthree\ and let $a$ be admissible; then every increasing sequence is a continuum subsequence for the Wick sector, every LIM state is $\Ginf$ ($=Z^a_{\infty,T}$ on $T_{\cF'}\cap\Dtube$), $\kappa^{\cK}=\kappa$ is independent of $\cK$, and Form~(B) of \cref{thm:s12:LIM} holds for $\Ginf$ with the full-sequence lattice clause.
\end{corollary}

\begin{proof}
Curvature sector: \cref{prop:s12:ID}(b) makes every increasing sequence a continuum subsequence with $\ZLIM{\cK}=Z_\infty$; (i)--(vi) are those of \cref{thm:s12:LIM}, and in (v) the local laws converge along the full sequence because the characteristic functionals now do (step ($\beta$) with the full sequence); (vii) is \cref{prop:s12:ID}(a). Wick sector: \cref{prop:s12:ID}(c); the clauses are Form~(B) of \cref{thm:s12:LIM} for $\Ginf$, with \OSt\ at amplitude $<H/2$ by the route of \cref{lem:s12:claim4prime}, which does not use \cref{hyp:s03:W1}, and at amplitude $<\sigma_W=\sigma_+/8$ by the route of \cref{thm:s08:OS2WT}; both hold ($H/2=\sigma_+/32<\sigma_+/8$: the route without the fixed-cube limits gives the smaller amplitude, and what it improves is the list of inputs); $\kappa^{\cK}=\kappa$ by the scheme clause.
\end{proof}

\Cref{hyp:s03:W1} is the scheme-$a$ analogue, with scalar sources, of the existence clause of King's theorem; it is not that theorem, which is false as printed with scalar sources in King's Wick ordering. With \Wtwo\ in addition one obtains Form~(A) for the full object $Z^a_\infty$ and, under \RW, its Euclidean clause; this is the optional extension of \cref{thm:s23:FormA}.

\begin{corollary}[King's convention]\label{cor:s12:KK}
Define continuum subsequences and LIM states of the K-family $a^K_{k,J}(z)=\avK{e^{\sum_jz_jF(f_j)}}$ (\cref{def:s02:Kcube}) by \cref{def:s12:contsub,def:s12:LIMstate} with $\avK\cdot$ in place of $\ava\cdot$.
\begin{enumerate}[label=\textup{(\alph*)},leftmargin=*,nosep]
\item \textup{(No hypothesis.)} Every increasing sequence of spacings has a continuum subsequence of the K-family; for every such $\cK$, family and $J\ge J_0(f)$ the fixed-cube limits $Z^{K,\cK}_J$ and the LIM state $Z^{K,\cK}_\infty=\lim_JZ^{K,\cK}_J$ exist, are entire and obey $|Z^{K,\cK}_J|,|Z^{K,\cK}_\infty|\le e^{\norm{g(z)}_2^2}$, with the other clauses of \cref{lem:s12:LIM0}(a) and \cref{lem:s12:LIM1}. Nothing beyond this is asserted for a K-family LIM state from the bounds alone.
\item \textup{(With \cref{hyp:s03:EC}.)} Every increasing sequence of spacings is a continuum subsequence of the K-family, every K-family LIM state is King's printed state $\ZKinf$, every LIM state of the comparison family is $\ZKinf$ as well (\cref{cor:s04:G2}(iii)), and the curvature-sector axioms of \cref{thm:s13:main} hold for it, with the lattice laws of \cref{thm:s09:S2K}(a) converging along the full sequence.
\end{enumerate}
With scalar sources, the corresponding statements on King's convention are supplied by \cref{thm:s13:TWK,prop:s13:WK}.
\end{corollary}

\begin{proof}
(a) The proof of \cref{lem:s12:LIM0}(a) uses only the $k$- and $J$-uniform local bound of \cref{lem:s05:C2} (Step~2), the second-order bound $\langle F(u)^2\rangle\le2\norm u^2_{2,k}$ of \cref{lem:s05:C1} (Step~3), and the current representation and Jensen's inequality (Step~4); that of \cref{lem:s12:LIM1} uses in addition the first inequality of \cref{lem:s02:cut}. All of these hold on King's convention: the bounds by \cref{thm:s05:CK} (the second-order bound is photon-only and transfers by the covariance identity of \cref{lem:s04:G}), the cut inequality by \cref{cor:s04:G1}. So the two lemmas hold for the K-family by the same proofs; the argument makes no use of \cref{hyp:s03:EC}. (b) By \cref{hyp:s03:EC}, \cref{cor:s04:G2}(ii)($\alpha$) gives, for every $J\ge J_0$ and family, locally uniform convergence of $a^K_{k,J}$ along the full sequence, so every increasing sequence is a continuum subsequence of the K-family, and by ($\beta$) every K-family LIM state is $\ZKinf$; \cref{cor:s04:G2}(iii) is the statement for the comparison family; the axioms are the curvature sector of \cref{thm:s13:main}, and the full-sequence lattice laws are \cref{thm:s09:S2K}(a).
\end{proof}

\begin{remark}[on the construction of the subsequences]\label{rem:s12:K9}
Every subsequence of spacings used in this paper is constructed in Step~2 of the proof of \cref{lem:s12:LIM0} from the $k$-uniform bounds of \cref{lem:s05:C2} and \eqref{eq:s11:W} and Montel's theorem, or is a further extraction from one so constructed, and the identification of its limits uses only the monotone cube limit of \cref{lem:s12:LIM1}; no subsequence statement is cited from \cite{K2}.
\end{remark}
